\section{Introduction}\label{sec:intro}

Drug discovery is a process of exceptional scientific and economic
complexity: the identification of a single clinical candidate
requires, on average, more than a decade and expenditures exceeding
one billion US dollars \cite{dimasi2016innovation}. A key bottleneck is the
experimental annotation of molecular properties---bioactivity,
toxicity, membrane permeability, and metabolic stability---for the
large chemical libraries assembled during early-stage screening.
Although high-throughput experimental techniques have expanded the
scale at which annotations can be acquired, the labelled data
available for most individual assays remains modest in comparison
with the requirements of modern supervised architectures \cite{walters2021applications}.

Machine learning has made substantial inroads into molecular property
prediction over the past decade. Quantitative structure-activity
relationship (QSAR) models based on molecular fingerprints, most
notably the Extended Connectivity Fingerprint (ECFP) family
\cite{wu2018moleculenet}, established strong baselines by encoding local atomic
environments as binary bit vectors amenable to conventional
classifiers. More recently, graph neural networks (GNNs) have
emerged as the dominant paradigm \cite{xiong2020pushing}, representing molecules
as attributed graphs and applying message-passing operations to
aggregate neighbourhood information into expressive molecular
representations. Nevertheless, the predictive power of these models
degrades under data-scarce conditions, since their high-dimensional
parameter spaces require large annotated corpora to avoid
overfitting \cite{guo2021fewshot}.

Data scarcity has motivated the development of few-shot and
meta-learning approaches. Altae-Tran et al.\ \cite{altaetran2017low} introduced
the iterative refinement LSTM combined with graph convolutions for
one-shot learning of molecular activities. Subsequent work proposed
model-agnostic meta-learning (MAML) \cite{vella2023fewshot}, prototypical
networks \cite{fang2023evidential}, and property-aware relational networks
\cite{wang2021property}, all of which exploit auxiliary tasks from related
assays to initialise models in parameter regions from which fast
adaptation with few examples is possible. Zhang et al.\ \cite{zhang2024molfescue}
further addressed the joint challenge of data scarcity and class
imbalance through a few-shot contrastive learning framework
(MolFeSCue) with a dynamic contrastive loss. Despite strong
empirical performance, these methods require either substantial
annotated auxiliary data or architectural choices whose optimality
for a given target cannot be guaranteed \textit{a priori}.

A conceptually attractive alternative is offered by
compression-based similarity measures rooted in algorithmic
information theory. The Normalized Compression Distance (NCD),
introduced by Cilibrasi and Vitányi \cite{cilibrasi2005clustering} as a computable
approximation to the theoretically optimal Normalized Information
Distance \cite{li2004similarity}, provides a parameter-free, domain-agnostic
measure of string similarity that has been applied successfully to
phylogenetic reconstruction, music classification, and text
clustering \cite{cilibrasi2005clustering,keogh2007compression}. The central insight is that if two
objects share informational content, a good compressor will
recognise this shared structure and compress their concatenation
more efficiently than the sum of their individual compressed
lengths. This principle requires no feature engineering, no
gradient computation, and no domain-specific knowledge.

This paradigm received renewed attention after Jiang et al.\
\cite{jiang2023lowresource} demonstrated that a gzip-based kNN classifier achieves
competitive performance with fine-tuned large language models on
several text classification benchmarks. Weinreich and Probst
\cite{weinreich2025learning} subsequently translated this finding to cheminformatics,
showing that NCD computed over SMILES strings achieves performance
comparable to ECFP-based kNN on BACE, BBBP, and aqueous solubility
tasks. These results suggest that \revmod{SMILES strings}
\cite{weininger1988smiles}---originally designed as a compact notation for
molecular graphs---encode sufficient structural information for
compression-based methods to recover biologically relevant
similarity patterns, consistent with the growing use of SMILES as
input for transformer-based molecular language models \cite{ross2022large}.

The present work extends the compression-based molecular
classification paradigm in four specific directions. First, SCoPE
aggregates signals from three complementary lossless compression
algorithms (gzip, zlib, bz2), each exploiting distinct statistical
regularities in SMILES. Second, instead of direct pairwise NCD
lookup for kNN classification, SCoPE constructs a centroid-based
geometric representation of each candidate class in the
dissimilarity matrix space, providing a more stable class signature
for small support sets. Third, SCoPE incorporates CDM \cite{keogh2007compression}
\revadd{and CLM} alongside NCD, capturing complementary aspects of
informational overlap. Fourth, SCoPE applies \revmod{the product of the
Euclidean and cosine distances} to the centroid--query comparison and
resolves the final prediction by majority voting \revmod{over the base
classifiers indexed by compressor--metric pairs}. The complete
implementation is publicly available at
\url{https://github.com/unlikeghost/ScOPE}.

The remainder of this paper is organised as follows.
Section~\ref{sec:methods} provides a full formal description of the
SCoPE pipeline. Section~\ref{sec:results} presents experimental
results, ablation studies, and a \emph{protocol-aware} comparative
analysis against ten related methods---an aspect that we treat
with particular care to avoid the misrepresentations that arise when
results from incompatible evaluation protocols are placed in the
same table without qualification. Section~\ref{sec:conclusion}
summarises the main findings.